\chapter{Captures des Interfaces Web et Mobile}
\minitoc
\newpage

\section{Introduction}
Ce chapitre est reserve a la presentation visuelle du projet GOSPORT. Il regroupe les emplacements prevus pour les captures des interfaces web et mobile. Les figures seront ajoutees apres l'execution finale du projet afin de montrer les ecrans reels obtenus.

\section{Interfaces Web}
Les captures web concernent principalement le tableau de bord Angular utilise par les administrateurs, les coachs et les nutritionnistes.

\subsection{Interface de connexion}
\begin{figure}[H]
\centering
\fbox{\parbox[c][7cm][c]{0.85\textwidth}{\centering Emplacement de la capture : interface de connexion web}}
\caption{Interface web de connexion}
\end{figure}

\subsection{Tableau de bord principal}
\begin{figure}[H]
\centering
\fbox{\parbox[c][7cm][c]{0.85\textwidth}{\centering Emplacement de la capture : tableau de bord web}}
\caption{Tableau de bord web GOSPORT}
\end{figure}

\subsection{Gestion des athlètes}
\begin{figure}[H]
\centering
\fbox{\parbox[c][7cm][c]{0.85\textwidth}{\centering Emplacement de la capture : liste et suivi des athlètes}}
\caption{Interface web de gestion des athlètes}
\end{figure}

\subsection{Création de programmes sportifs}
\begin{figure}[H]
\centering
\fbox{\parbox[c][7cm][c]{0.85\textwidth}{\centering Emplacement de la capture : creation ou edition d'un programme}}
\caption{Interface web de création de programme}
\end{figure}

\subsection{Module nutritionnel}
\begin{figure}[H]
\centering
\fbox{\parbox[c][7cm][c]{0.85\textwidth}{\centering Emplacement de la capture : dashboard nutrition ou diet builder}}
\caption{Interface web du module nutritionnel}
\end{figure}

\subsection{Messagerie et sessions}
\begin{figure}[H]
\centering
\fbox{\parbox[c][7cm][c]{0.85\textwidth}{\centering Emplacement de la capture : messagerie, notifications ou calendrier}}
\caption{Interface web de communication et sessions}
\end{figure}

\section{Interfaces Mobile}
Les captures mobile concernent l'application Flutter destinee principalement aux athlètes. Elles illustrent les parcours de consultation, d'execution et de suivi quotidien.

\subsection{Connexion mobile}
\begin{figure}[H]
\centering
\fbox{\parbox[c][10cm][c]{0.55\textwidth}{\centering Emplacement de la capture : ecran de connexion mobile}}
\caption{Interface mobile de connexion}
\end{figure}

\subsection{Accueil mobile}
\begin{figure}[H]
\centering
\fbox{\parbox[c][10cm][c]{0.55\textwidth}{\centering Emplacement de la capture : accueil ou dashboard mobile}}
\caption{Tableau de bord mobile}
\end{figure}

\subsection{Programme et séance}
\begin{figure}[H]
\centering
\fbox{\parbox[c][10cm][c]{0.55\textwidth}{\centering Emplacement de la capture : programme, detail ou workout player}}
\caption{Interface mobile de suivi d'entraînement}
\end{figure}

\subsection{Nutrition mobile}
\begin{figure}[H]
\centering
\fbox{\parbox[c][10cm][c]{0.55\textwidth}{\centering Emplacement de la capture : suivi nutritionnel mobile}}
\caption{Interface mobile de nutrition}
\end{figure}

\subsection{Profil et notifications}
\begin{figure}[H]
\centering
\fbox{\parbox[c][10cm][c]{0.55\textwidth}{\centering Emplacement de la capture : profil, notifications ou calendrier}}
\caption{Interface mobile de profil et notifications}
\end{figure}

\section{Conclusion}
Les emplacements ci-dessus seront remplaces par les captures finales de l'application. Cette organisation permet de documenter separement les interfaces web et mobile tout en conservant une numerotation claire dans la table des figures.
